\documentclass[11pt]{article}

% ===================== PACKAGES =====================
\usepackage[a4paper,margin=1in]{geometry}
\usepackage{amsmath,amssymb}
\usepackage{enumitem}
\usepackage{fancyhdr}
\usepackage{xcolor}

% ===================== HEADER =====================
\pagestyle{fancy}
\fancyhf{}
\lhead{CSE 400: Fundamentals of Probability in Computing}
\rhead{Tutorial 1 – Lecture Scribe}
\cfoot{\thepage}

% ===================== TITLE =====================
\title{
\normalsize School of Engineering and Applied Science (SEAS), Ahmedabad University \\
\vspace{0.2cm}
\textbf{CSE 400: Fundamentals of Probability in Computing}\\
\Large Tutorial Review Session Scribe
}
\author{}
\date{}

\begin{document}

\maketitle

\vspace{-1cm}

\begin{center}
\begin{tabular}{ll}
\textbf{Name:} & Khushi \\[1.5ex]
\textbf{Enrollment No:} & AU2440131 \\[1.5ex]
\textbf{Course:} & CSE 400 \\[1.5ex]
\textbf{Topic:} & Tutorial 1 Review
\end{tabular}
\end{center}

\hrule
\vspace{0.5cm}

% ===================================================
\section*{Objective}

This scribe serves as a structured revision checkpoint and academic reference for 
Tutorial Session 1 of \textbf{CSE 400: Fundamentals of Probability in Computing}. 
The document reviews key probability distributions and counting principles used in solving the tutorial problems.

\bigskip

% ===================================================
\section*{Core Formulas}

\textbf{1. Poisson Distribution}

\[
P(X=k) = \frac{e^{-\lambda}\lambda^k}{k!}
\]

\textbf{2. Normal Distribution PDF}

\[
f_X(x)=\frac{1}{\sqrt{2\pi\sigma^2}}
\exp\left(-\frac{(x-\mu)^2}{2\sigma^2}\right)
\]

\textbf{3. Conditional Probability}

\[
P(A|B) = \frac{P(A \cap B)}{P(B)}
\]

\textbf{4. Binomial Distribution}

\[
P(X=k) = \binom{n}{k}p^k(1-p)^{n-k}
\]

\bigskip

\hrule
\vspace{0.5cm}

% ===================================================
\section*{Tutorial Problems and Solutions}

\begin{enumerate}[label=\textbf{Question \arabic*:}]

%----------------------------------------------------
\item \textbf{Food Festival Dish Division}

We divide 20 distinct dishes into four groups of 5.

\textbf{Total ways}

\[
\frac{20!}{(5!)^4 4!}
\]

\textbf{Probability each platter has one cuisine}

If each cuisine already has 5 dishes, there is only one favourable arrangement.

\[
P = \frac{1}{\frac{20!}{(5!)^4 4!}}
\]

\newpage

%----------------------------------------------------
\item \textbf{Bus Waiting Time}

Passenger arrival is uniform between 0 and 30 minutes.

\textbf{(a) Wait < 5 minutes}

Favourable interval = 10 minutes

\[
P = \frac{10}{30} = \frac{1}{3}
\]

\textbf{(b) Wait > 10 minutes}

Favourable interval = 10 minutes

\[
P = \frac{10}{30} = \frac{1}{3}
\]

\newpage

%----------------------------------------------------
\item \textbf{Conditional Probability}

Given

\[
P(A)=0.15, \quad P(B)=0.25, \quad P(A\cap B)=0.08
\]

\[
P(A|B)=\frac{0.08}{0.25}
\]

\[
P(A|B)=0.32
\]

\newpage

%----------------------------------------------------
\item \textbf{Typographical Errors}

\[
\lambda = 0.2
\]

\textbf{(a) Probability of 0 errors}

\[
P(X=0)=e^{-0.2}
\]

\[
P(X=0)\approx0.8187
\]

\textbf{(b) Probability of at least 2 errors}

\[
P(X\ge2)=1-P(0)-P(1)
\]

\[
P(1)=0.2e^{-0.2}
\]

\newpage

%----------------------------------------------------
\item \textbf{Airplane Crashes}

\[
\lambda =3.5
\]

\textbf{(a) At least 2 crashes}

\[
P(X\ge2)=1-P(0)-P(1)
\]

\textbf{(b) At most 1 crash}

\[
P(X\le1)=P(0)+P(1)
\]

\newpage

%----------------------------------------------------
\item \textbf{Highway Accidents}

\[
X\sim Poisson(3)
\]

\textbf{(a)}

\[
P(X\ge3)=1-P(0)-P(1)-P(2)
\]

\textbf{(b)}

\[
P(X\ge3|X\ge1)=\frac{P(X\ge3)}{P(X\ge1)}
\]

\newpage

%----------------------------------------------------
\item \textbf{Normal Distribution}

Given

\[
f_X(x)=\frac{1}{\sqrt{8\pi}} e^{-\frac{(x+3)^2}{8}}
\]

Comparing with the standard normal form:

\[
\mu=-3,\quad \sigma=2
\]

\textbf{Example}

\[
Z=\frac{X-\mu}{\sigma}
\]

\newpage

%----------------------------------------------------
\item \textbf{Jury Decision}

\[
X\sim Binomial(12,\theta)
\]

Conviction occurs if at least 8 jurors vote guilty.

\[
P(\text{correct})=
\sum_{k=8}^{12}\binom{12}{k}\theta^k(1-\theta)^{12-k}
\]

\newpage

%----------------------------------------------------
\item \textbf{PMF Identification}

Given

\[
p(i)=\frac{c\lambda^i}{i!}
\]

Using normalization:

\[
\sum_{i=0}^{\infty}p(i)=1
\]

\[
c=e^{-\lambda}
\]

Thus

\[
p(i)=\frac{e^{-\lambda}\lambda^i}{i!}
\]

which is a Poisson distribution.

\newpage

%----------------------------------------------------
\item \textbf{Communication System Reliability}

Each component works with probability $p$.

\[
X\sim Binomial(n,p)
\]

System works if at least half of components function.

For 3 components:

\[
P_3=P(X\ge2)
\]

For 5 components:

\[
P_5=P(X\ge3)
\]

The larger system is better when

\[
p>\frac{1}{2}
\]

\end{enumerate}

\bigskip

\hrule
\vspace{0.3cm}

\begin{center}
\small \textit{End of Tutorial Scribe}
\end{center}

\end{document}
